\documentclass[12pt,a4paper]{article}

% ============================================================================
% PACKAGES
% ============================================================================
\usepackage[utf8]{inputenc}
\usepackage[T1]{fontenc}
\usepackage{amsmath,amssymb,amsfonts}
\usepackage{graphicx}
\usepackage{booktabs}
\usepackage{multirow}
\usepackage{array}
\usepackage{hyperref}
\usepackage{xcolor}
\usepackage{algorithm}
\usepackage{algpseudocode}
\usepackage{subcaption}
\usepackage[margin=1in]{geometry}
\usepackage{natbib}
\usepackage{siunitx}

% ============================================================================
% CUSTOM COMMANDS
% ============================================================================
\newcommand{\wss}{\tau_w}
\newcommand{\pinn}{PINN}
\newcommand{\cfd}{CFD}
\newcommand{\nrmse}{NRMSE}
\newcommand{\rmse}{RMSE}

% ============================================================================
% TITLE AND AUTHORS
% ============================================================================
\title{Computational Investigation of Blood Flow in Saphenous Vein Grafts, Left and Right Coronary Arteries: Healthy and Diseased Cases}

\author{
    [Author Names] \\
    \textit{[Affiliation]} \\
    \textit{[Email]}
}

\date{\today}

% ============================================================================
% DOCUMENT
% ============================================================================
\begin{document}

\maketitle

\begin{abstract}
Wall shear stress (WSS) is a critical hemodynamic parameter for predicting cardiovascular outcomes in patients with coronary artery disease (CAD) and those who have undergone coronary artery bypass grafting (CABG) with saphenous vein grafts (SVG). This study presents a comprehensive computational investigation of blood flow in coronary arteries and bypass grafts across healthy controls, diseased cases, and post-CABG patients. We employ computational fluid dynamics (CFD) simulations using both Newtonian and Carreau-Yasuda non-Newtonian blood models to characterize WSS distributions. Additionally, we develop physics-informed neural networks (PINNs) as surrogate models that learn to predict WSS while satisfying the incompressible Navier-Stokes equations, enabling real-time hemodynamic assessment. We demonstrate that direct WSS prediction significantly outperforms approaches that compute WSS from velocity gradients, due to fundamental conflicts with the no-slip boundary condition. Our FourierPINN architecture, trained on CFD data from 9 patient cases (2 healthy, 1 diseased, 6 SVG), achieves a mean normalized root-mean-square error (NRMSE) of 0.61\% and mean $R^2$ of 0.967 compared to CFD ground truth. The trained models predict WSS distributions across 145,000--388,000 wall points in under 1 second, representing a speedup of 3--4 orders of magnitude over traditional CFD simulations. This approach enables rapid hemodynamic evaluation for clinical decision support in CABG surgery planning and post-operative monitoring.
\end{abstract}

\textbf{Keywords:} Wall Shear Stress, Coronary Artery Disease, Saphenous Vein Graft, Computational Fluid Dynamics, Physics-Informed Neural Networks, CABG, Deep Learning

% ============================================================================
% SECTION 1: INTRODUCTION
% ============================================================================
\section{Introduction}

Coronary artery disease (CAD) remains the leading cause of mortality worldwide, accounting for approximately 17.9 million deaths annually. Coronary artery bypass grafting (CABG) surgery is a primary intervention for severe multi-vessel CAD, with saphenous vein grafts (SVG) being the most commonly used conduits \cite{cabg_outcomes}. However, SVG failure rates remain significant, with 10--20\% of grafts occluding within the first year and up to 50\% failing within 10 years.

The long-term patency of saphenous vein grafts is significantly influenced by local hemodynamic conditions, particularly wall shear stress (WSS) \cite{wss_graft_failure}. WSS represents the tangential force per unit area exerted by flowing blood on the vessel wall, playing a crucial role in endothelial cell function and vascular remodeling. Low WSS ($<$ 4 dyn/cm$^2$ or 0.4 Pa) promotes atherosclerotic plaque formation through increased endothelial permeability and inflammatory cell adhesion, while abnormally high WSS ($>$ 40 dyn/cm$^2$ or 4 Pa) can cause endothelial damage and intimal hyperplasia \cite{wss_thresholds}.

Computational fluid dynamics (CFD) has become an essential tool for non-invasive WSS assessment, enabling patient-specific hemodynamic analysis from medical imaging data \cite{cfd_coronary}. CFD provides detailed spatial information about blood flow patterns that cannot be obtained through clinical imaging alone. However, CFD simulations face significant practical limitations that hinder clinical adoption. High-fidelity simulations require hours to days of computation on specialized hardware, while mesh generation, boundary condition specification, and result interpretation demand specialized expertise. Furthermore, design optimization or uncertainty quantification requires multiple expensive simulations, making iterative workflows impractical. These time and expertise requirements ultimately limit routine clinical application of CFD for surgical planning.

Physics-informed neural networks (PINNs) offer a promising approach to overcome these limitations by embedding physical laws directly into the neural network training process \cite{raissi2019physics}. Unlike purely data-driven machine learning methods, PINNs enforce the governing equations (e.g., Navier-Stokes) as soft constraints during training, enabling accurate predictions while ensuring physical consistency.

In this work, we present a comprehensive computational investigation of blood flow in coronary arteries and saphenous vein grafts that combines traditional CFD analysis with physics-informed machine learning. We first conduct CFD simulations comparing Newtonian and Carreau-Yasuda non-Newtonian blood models across healthy, diseased, and post-CABG cases. Building on these simulations, we develop a FourierPINN architecture with Fourier feature encoding that overcomes spectral bias for learning sharp WSS gradients. We validate the PINN surrogate model across 9 patient cases spanning diverse vascular anatomies and demonstrate real-time inference capability enabling clinical translation.

% ============================================================================
% SECTION 2: METHODS
% ============================================================================
\section{Methods}

\subsection{Governing Equations and Numerical Modeling}

Blood, composed of cells suspended in plasma, demonstrates non-Newtonian behavior, particularly notable in vessels with diameters smaller than 1 mm. Under certain pathological and transient conditions, such as artery dilation and bifurcations, shear rates may decrease below 100 s$^{-1}$, emphasizing the potential importance of non-Newtonian processes. Conversely, in larger blood vessels, where shear rates typically exceed 100 s$^{-1}$ and with diameters exceeding 0.3 mm, blood can be accurately represented as a Newtonian fluid.

In our current study, we compare the Newtonian model with the non-Newtonian Carreau-Yasuda model. The models incorporate a rigid artery wall and account for blood density of $\rho = 1050$ kg/m$^3$ and viscosity of $\mu = 0.0035$ Pa$\cdot$s for the Newtonian model. The incompressible Navier-Stokes equations characterize the flow behavior:

\begin{equation}
    \nabla \cdot \mathbf{v} = 0
    \label{eq:continuity}
\end{equation}

\begin{equation}
    \rho(\mathbf{v} \cdot \nabla)\mathbf{v} = -\nabla p + \mu \nabla^2 \mathbf{v}
    \label{eq:navier_stokes}
\end{equation}

\noindent where $\rho$ denotes density and $\mu$ represents dynamic viscosity.

For the non-Newtonian Carreau-Yasuda model, the viscosity $\mu(\dot{\gamma})$ is given by:

\begin{equation}
    \mu(\dot{\gamma}) = \mu_\infty + (\mu_0 - \mu_\infty) \left(1 + (\lambda \dot{\gamma})^a\right)^{\frac{n-1}{a}}
    \label{eq:carreau_yasuda}
\end{equation}

\noindent where $\dot{\gamma}$ is the scalar shear rate, and the parameters are: $\mu_\infty = 0.0035$ Pa$\cdot$s, $\mu_0 = 0.16$ Pa$\cdot$s, $\lambda = 8.2$ s, $n = 0.2128$, and $a = 0.64$.

A uniform velocity of 10 cm/s is applied at the inlet for all models, ensuring laminar flow. Consistent with published data, 4\% of the inlet flow is directed to the coronary circulation. Rigid arterial walls with no-slip conditions are assumed. The WSS is calculated using:

\begin{equation}
    \tau_w = \mu \left. \frac{\partial v}{\partial n} \right|_{n=0}
    \label{eq:wss}
\end{equation}

\noindent where $\frac{\partial v}{\partial n}$ is the velocity gradient normal to the wall.

\subsection{CFD Simulation}

The governing equations are solved using ANSYS Fluent with the finite volume method. The pressure-based solver is employed with residual convergence threshold of $10^{-3}$ to $10^{-4}$. Following a mesh independence study, optimal mesh size ranges between 0.3 mm and 0.4 mm.

\subsection{Physics-Informed Neural Network Formulation}

To enable rapid WSS prediction without repeated CFD simulations, we developed physics-informed neural networks (PINNs) as surrogate models. PINNs approximate the solution to PDEs by training a neural network $\mathcal{N}_\theta: \mathbb{R}^3 \rightarrow \mathbb{R}^5$ that maps spatial coordinates $(x, y, z)$ to hemodynamic quantities $(u, v, w, p, \tau_w)$.

The training objective minimizes a composite loss function:
\begin{equation}
    \mathcal{L} = \lambda_{\text{wss}} \mathcal{L}_{\text{wss}} + \lambda_{\text{vel}} \mathcal{L}_{\text{vel}} + \lambda_{\text{NS}} \mathcal{L}_{\text{NS}} + \lambda_{\text{cont}} \mathcal{L}_{\text{cont}}
\end{equation}

\subsubsection{Data Loss Terms}
The WSS data loss enforces agreement with CFD results at wall points:
\begin{equation}
    \mathcal{L}_{\text{wss}} = \frac{1}{N_{\text{wall}}} \sum_{i=1}^{N_{\text{wall}}} |\tau_{w,i}^{\text{pred}} - \tau_{w,i}^{\text{CFD}}|^2
\end{equation}

The velocity data loss ensures accurate velocity field reconstruction:
\begin{equation}
    \mathcal{L}_{\text{vel}} = \frac{1}{N} \sum_{i=1}^{N} \|\mathbf{u}_i^{\text{pred}} - \mathbf{u}_i^{\text{CFD}}\|^2
\end{equation}

\subsubsection{Navier-Stokes Residual}
The physics loss enforces the incompressible Navier-Stokes momentum equations:
\begin{equation}
    \rho(\mathbf{u} \cdot \nabla)\mathbf{u} = -\nabla p + \mu \nabla^2 \mathbf{u}
\end{equation}

The residual for each momentum component is:
\begin{align}
    f_u &= \rho\left(u\frac{\partial u}{\partial x} + v\frac{\partial u}{\partial y} + w\frac{\partial u}{\partial z}\right) + \frac{\partial p}{\partial x} - \mu\left(\frac{\partial^2 u}{\partial x^2} + \frac{\partial^2 u}{\partial y^2} + \frac{\partial^2 u}{\partial z^2}\right) \\
    f_v &= \rho\left(u\frac{\partial v}{\partial x} + v\frac{\partial v}{\partial y} + w\frac{\partial v}{\partial z}\right) + \frac{\partial p}{\partial y} - \mu\left(\frac{\partial^2 v}{\partial x^2} + \frac{\partial^2 v}{\partial y^2} + \frac{\partial^2 v}{\partial z^2}\right) \\
    f_w &= \rho\left(u\frac{\partial w}{\partial x} + v\frac{\partial w}{\partial y} + w\frac{\partial w}{\partial z}\right) + \frac{\partial p}{\partial z} - \mu\left(\frac{\partial^2 w}{\partial x^2} + \frac{\partial^2 w}{\partial y^2} + \frac{\partial^2 w}{\partial z^2}\right)
\end{align}

The Navier-Stokes loss is:
\begin{equation}
    \mathcal{L}_{\text{NS}} = \frac{1}{N_c} \sum_{i=1}^{N_c} \left( f_u^2 + f_v^2 + f_w^2 \right)
\end{equation}
where $N_c$ is the number of collocation points sampled from the computational domain.

\subsubsection{Continuity Residual}
Mass conservation for incompressible flow:
\begin{equation}
    \mathcal{L}_{\text{cont}} = \frac{1}{N_c} \sum_{i=1}^{N_c} \left( \frac{\partial u}{\partial x} + \frac{\partial v}{\partial y} + \frac{\partial w}{\partial z} \right)^2
\end{equation}

\subsubsection{WSS Physics Constraint}
Wall shear stress is related to velocity gradients at the wall:
\begin{equation}
    \wss = \mu \left| \frac{\partial \mathbf{u}_{\text{tangent}}}{\partial n} \right|_{n=0}
\end{equation}

For no-slip walls where $\mathbf{u} = 0$, this simplifies to:
\begin{equation}
    \wss = \mu \sqrt{\left(\frac{\partial u}{\partial n}\right)^2 + \left(\frac{\partial v}{\partial n}\right)^2 + \left(\frac{\partial w}{\partial n}\right)^2}
\end{equation}

The WSS physics loss enforces consistency between predicted and gradient-computed WSS:
\begin{equation}
    \mathcal{L}_{\text{wss}} = \frac{1}{N_{\text{wall}}} \sum_{i=1}^{N_{\text{wall}}} \left| \wss_i^{\text{pred}} - \wss_i^{\text{computed}} \right|^2
\end{equation}

\subsection{FourierPINN Architecture}

Standard neural networks suffer from \textit{spectral bias}---the tendency to learn low-frequency functions before high-frequency components \cite{spectral_bias}. This is problematic for WSS prediction, where sharp gradients occur at bifurcations and stenoses.

We employ Fourier feature encoding \cite{tancik2020fourier} to overcome this limitation:
\begin{equation}
    \gamma(\mathbf{x}) = \left[ \mathbf{x}, \sin(2\pi \mathbf{B} \mathbf{x}), \cos(2\pi \mathbf{B} \mathbf{x}) \right]
\end{equation}
where $\mathbf{B} \in \mathbb{R}^{m \times 3}$ is a random Gaussian matrix with $m = 64$ frequencies and scale $\sigma = 10$.

The FourierPINN architecture takes 3D coordinates $(x, y, z)$ as input and maps them through the Fourier encoding to a $(3 + 2 \times 64) = 131$ dimensional representation. This encoded input passes through a trunk network consisting of 8 fully-connected layers with 128 units each, using SiLU activation functions. The network terminates with separate linear output heads for each predicted quantity: the three velocity components $(u, v, w)$, pressure $p$, and wall shear stress $\wss$. The total parameter count is approximately 133,000, making the model lightweight for deployment.

\subsection{Training Procedure}

\subsubsection{Data Preprocessing}
Coordinates are normalized to $[0, 1]$ using min-max scaling:
\begin{equation}
    \tilde{x} = \frac{x - x_{\min}}{x_{\max} - x_{\min}}
\end{equation}

Gradients in the physics loss require chain rule correction:
\begin{equation}
    \frac{\partial u}{\partial x} = \frac{\partial u}{\partial \tilde{x}} \cdot \frac{1}{x_{\max} - x_{\min}}
\end{equation}

Velocities are normalized to $[-1, 1]$ and WSS to $[0, 1]$.

\subsubsection{Collocation Point Sampling}
Collocation points for physics loss evaluation are sampled from the CFD mesh rather than uniformly from the bounding box. This is critical because coronary arteries are thin tubular structures (diameter 3--5 mm) within a much larger bounding volume. Random uniform sampling would generate points outside the vessel where the Navier-Stokes equations do not apply.

Interior (streamline) points are weighted higher than wall points since they provide more informative physics residuals (non-zero velocity gradients).

\subsubsection{Optimization}
The network is trained using the AdamW optimizer with a weight decay of $10^{-5}$ to prevent overfitting. The learning rate is initialized at $2 \times 10^{-4}$ and follows a cosine annealing schedule, decaying to $10^{-6}$ over the training period. Each training batch contains 4,096 data points, supplemented with 2,048 collocation points for physics loss evaluation. To prevent overfitting and reduce unnecessary computation, early stopping is employed with a patience of 20 epochs, monitoring the total loss. Gradient clipping with a maximum norm of 1.0 ensures stable optimization, particularly important given the second-order derivatives in the Navier-Stokes residuals.

\subsection{CFD Ground Truth Generation}

CFD simulations were performed using ANSYS Fluent to generate training data for the PINN models. Blood was modeled as a Newtonian fluid with density $\rho = 1050$ kg/m$^3$ and dynamic viscosity $\mu = 0.0035$ Pa$\cdot$s. A uniform velocity of 10 cm/s was applied at the inlet, with pressure outlets configured to direct 4\% of the inlet flow to the coronary circulation, consistent with physiological data. Rigid arterial walls with no-slip boundary conditions were assumed. Following a mesh independence study, the computational mesh was refined to element sizes between 0.3--0.4 mm. Simulations were considered converged when residuals fell below $10^{-3}$ to $10^{-4}$.

\subsection{Patient Cohort}

The study includes 9 patient cases from multiple sources representing diverse clinical presentations. Healthy controls (H-09 and H-12) with normal coronary anatomy were obtained from the Vascular Model Repository. A diseased case (D-10) with stenosed coronary arteries was sourced from the Automated Segmentation of Coronary Arteries (ASOCA) dataset. The remaining six cases (0073, 0148, 0149, 0150, 0156, and ND2) represent post-CABG patients with saphenous vein grafts, derived from clinical CT angiography.

For CAD patients without CABG, both right and left coronary artery measurements are reported. For CABG patients, coronary artery measurements were reported as long as they were functional either spontaneously or with bypassed grafts.

Table \ref{tab:patient_data} summarizes the dataset characteristics.

\begin{table}[htbp]
\centering
\caption{Patient dataset characteristics from CFD simulations}
\label{tab:patient_data}
\begin{tabular}{lccccc}
\toprule
\textbf{Patient} & \textbf{Category} & \textbf{Vessels} & \textbf{Wall Points} & \textbf{WSS Range (Pa)} \\
\midrule
H-12 & Healthy & Aorta, LCA & 254,492 & 0--94 \\
H-09 & Healthy & Aorta, RCA & 145,432 & 0--99 \\
D-10 & Diseased & Aorta, LCA, RCA & 267,767 & 0--140 \\
0073 & Mixed & Aorta, LCA, RCA & 321,448 & 0--13 \\
0149 & SVG & Aorta, G1, G2, G3 & 312,911 & 0--67 \\
0156 & SVG & Aorta, G2, G3 & 387,696 & 0--35 \\
0148 & SVG & Aorta, G2 & 272,684 & 0--31 \\
0150 & SVG & Aorta, G3 & 213,961 & 0--168 \\
ND2 & Unknown & Aorta, LCA & 149,972 & 0--1015 \\
\bottomrule
\end{tabular}
\end{table}

Elevated WSS was noted at the outlets of some coronary arteries and grafts. To enhance accuracy, we clipped 1--2 cm from outlet regions to mitigate boundary effects and focus on physiologically representative flow conditions.

\subsection{Evaluation Metrics}

Model performance was assessed using four complementary metrics. Root mean square error (RMSE) quantifies the average magnitude of prediction errors in physical units (Pa), computed as $\sqrt{\frac{1}{N}\sum(y_i - \hat{y}_i)^2}$. Mean absolute error (MAE), calculated as $\frac{1}{N}\sum|y_i - \hat{y}_i|$, provides a robust measure less sensitive to outliers. To enable comparison across patients with different WSS ranges, we computed the normalized RMSE (NRMSE) as the ratio of RMSE to the data range $(y_{\max} - y_{\min})$. Finally, the coefficient of determination ($R^2 = 1 - \frac{\sum(y_i - \hat{y}_i)^2}{\sum(y_i - \bar{y})^2}$) quantifies how well predictions explain variance in the CFD ground truth.

\subsection{Statistical Analysis}

Statistical analysis was performed to evaluate WSS parameters across vessel types. The computed variables included area-weighted average WSS, average WSS at bifurcations, area fraction where WSS exceeds 40 dyn/cm$^2$, and interquartile range (IQR) of WSS. The Mann-Whitney U test was employed to compare parameters between healthy controls and CAD patients, with significance threshold $p < 0.05$.

% ============================================================================
% SECTION 3: RESULTS
% ============================================================================
\section{Results}

\subsection{CFD Simulation Results}

[This section presents CFD results comparing Newtonian and Carreau-Yasuda models, WSS distributions in healthy vs diseased cases, and statistical analysis of WSS parameters across patient groups. Include WSS contour plots and statistical comparisons.]

\subsection{PINN Training and Convergence}

All PINN models converged within 124--235 epochs with early stopping, well before the maximum of 5,000 epochs. Training times ranged from 38 minutes (H-09) to 187 minutes (0156), with an average of approximately 90 minutes per patient on an NVIDIA GeForce RTX 5060 Ti GPU.

\subsection{PINN Prediction Accuracy}

Table \ref{tab:results} presents the quantitative evaluation of PINN predictions against CFD ground truth for all patients.

\begin{table}[htbp]
\centering
\caption{FourierPINN WSS prediction performance compared to CFD ground truth}
\label{tab:results}
\begin{tabular}{lcccccc}
\toprule
\textbf{Patient} & \textbf{Category} & \textbf{RMSE (Pa)} & \textbf{MAE (Pa)} & \textbf{NRMSE (\%)} & \textbf{$R^2$} & \textbf{Epochs} \\
\midrule
H-12 & Healthy & 0.410 & 0.234 & 0.43 & 0.977 & 124 \\
H-09 & Healthy & 0.556 & 0.267 & 0.56 & 0.961 & 132 \\
D-10 & Diseased & 0.726 & 0.411 & 0.52 & 0.967 & 212 \\
0073 & Mixed & 0.129 & 0.103 & 0.96 & 0.965 & 195 \\
0149 & SVG & 0.275 & 0.111 & 0.41 & 0.953 & 138 \\
0156 & SVG & 0.279 & 0.146 & 0.80 & 0.982 & 235 \\
\midrule
\textbf{Mean} & -- & \textbf{0.396} & \textbf{0.212} & \textbf{0.61} & \textbf{0.967} & \textbf{173} \\
\textbf{Std} & -- & 0.200 & 0.111 & 0.21 & 0.011 & 45 \\
\bottomrule
\end{tabular}
\end{table}

The results demonstrate excellent agreement between PINN predictions and CFD ground truth across all patient categories. The mean NRMSE of 0.61\% indicates that prediction errors are less than 1\% of the WSS range, while the mean $R^2$ of 0.967 confirms strong correlation between predicted and true values. Importantly, RMSE values remained below 1 Pa for all patients, which is clinically insignificant given typical WSS ranges of 0--100+ Pa in coronary vessels. Performance was notably consistent across healthy, diseased, and SVG categories, suggesting the architecture generalizes well to diverse vascular geometries.

\subsection{Qualitative Comparison}

\begin{figure}[htbp]
\centering
% \includegraphics[width=\textwidth]{figures/wss_convergence_comparison.png}
\caption{WSS loss convergence comparison between direct prediction and velocity-derived approaches. Direct WSS prediction (blue) shows steady decrease from 0.04 to below 0.001, while velocity-derived WSS (orange) remains stagnant at approximately 0.0012 throughout training. This demonstrates the fundamental limitation of computing WSS from velocity gradients at no-slip wall boundaries.}
\label{fig:wss_convergence_comparison}
\end{figure}

Figure \ref{fig:wss_comparison} shows representative WSS distributions comparing CFD ground truth with PINN predictions. The PINN accurately captures low WSS regions in the aortic root and large vessel segments, as well as elevated WSS at bifurcations and stenotic regions where flow acceleration occurs. Spatial gradients along vessel walls are faithfully reproduced, and the complex WSS patterns characteristic of saphenous vein grafts are well represented. These qualitative observations confirm that the FourierPINN architecture successfully learns both the overall WSS distribution and the fine spatial details critical for clinical interpretation.

\subsection{Computational Performance}

Table \ref{tab:timing} compares computational requirements between CFD and PINN approaches.

\begin{table}[htbp]
\centering
\caption{Computational performance comparison}
\label{tab:timing}
\begin{tabular}{lcc}
\toprule
\textbf{Metric} & \textbf{CFD} & \textbf{PINN} \\
\midrule
Mesh generation & 30--60 min & N/A \\
Simulation time & 2--8 hours & N/A \\
Training time (one-time) & N/A & 2--3 hours \\
Inference time & N/A & $<$ 1 second \\
Hardware & 32+ CPU cores & Single GPU \\
\midrule
\textbf{Total per case} & \textbf{3--9 hours} & \textbf{$<$ 1 second}$^*$ \\
\bottomrule
\end{tabular}
\begin{tablenotes}
\small
\item $^*$After one-time training; training required once per patient or can be fine-tuned
\end{tablenotes}
\end{table}

The PINN achieves 3--4 orders of magnitude speedup at inference time, fundamentally changing what is computationally feasible in clinical settings. Surgeons can now perform real-time WSS visualization during surgical planning, interactively exploring hemodynamic parameters across different graft configurations. The rapid inference also enables uncertainty quantification through ensemble predictions, providing confidence intervals alongside point estimates. Most significantly, this computational efficiency makes integration into clinical decision support systems practical, bridging the gap between research-grade CFD analysis and bedside utility.

% ============================================================================
% SECTION 4: DISCUSSION
% ============================================================================
\section{Discussion}

\subsection{CFD Model Comparison}

[Discussion of Newtonian vs Carreau-Yasuda model differences. In larger coronary arteries where shear rates exceed 100 s$^{-1}$, the Newtonian assumption provides accurate results. Non-Newtonian effects become more pronounced in low-flow regions and at bifurcations where shear rates decrease.]

\subsection{Clinical Implications of WSS Findings}

[Discussion of WSS patterns in healthy vs diseased cases, clinical significance of elevated/reduced WSS regions, and implications for graft patency.]

\subsection{Direct WSS Prediction vs. Velocity-Derived WSS}

A fundamental design choice in physics-informed neural networks for WSS prediction is whether to predict WSS directly as a network output, or to compute WSS from the predicted velocity field using the physical relationship:
\begin{equation}
    \wss = \mu \sqrt{\left(\frac{\partial u}{\partial n}\right)^2 + \left(\frac{\partial v}{\partial n}\right)^2 + \left(\frac{\partial w}{\partial n}\right)^2} \bigg|_{\text{wall}}
    \label{eq:wss_from_velocity}
\end{equation}

We investigated both approaches and found that direct WSS prediction significantly outperforms the velocity-derived approach. This section explains the underlying reasons and provides guidance for future PINN implementations in hemodynamics.

\subsubsection{The No-Slip Boundary Condition Conflict}

The velocity-derived WSS approach faces a fundamental mathematical conflict at wall boundaries. The no-slip boundary condition requires:
\begin{equation}
    \mathbf{u}|_{\text{wall}} = 0
\end{equation}

This means velocity is identically zero at the wall surface, creating a singularity for gradient computation. While the velocity gradient $\nabla \mathbf{u}$ should be non-zero at the wall (since velocity transitions from zero to the free-stream value), the neural network must learn this gradient from training data where all wall points have $\mathbf{u} = 0$.

In practice, we observed that the velocity-derived WSS loss remained constant at approximately 0.0012 throughout training, while the direct WSS prediction loss decreased steadily from 0.04 to below 0.001 (Figure \ref{fig:wss_convergence_comparison}). This indicates that the network cannot effectively learn the near-wall velocity gradients from wall-point supervision alone.

\subsubsection{Attempted Remediation Strategies}

We explored several strategies to improve velocity-derived WSS computation, though none proved successful.

Rather than computing gradients exactly at wall points where $\mathbf{u} = 0$, we attempted to offset evaluation points into the domain along the inward surface normal, using $\mathbf{x}_{\text{offset}} = \mathbf{x}_{\text{wall}} + \epsilon \cdot \mathbf{n}$ where $\epsilon = 0.02$ in normalized coordinates and $\mathbf{n}$ is the inward-pointing normal. This approach allows gradient evaluation at points with non-zero velocity. However, it introduced extrapolation errors and did not resolve the fundamental learning difficulty.

We also ensured proper gradient flow during backpropagation by avoiding detach operations on intermediate WSS computations. While necessary for optimization, this modification did not address the underlying data limitation.

Increasing the velocity loss weight relative to the physics loss was another strategy we evaluated, with the goal of emphasizing accurate velocity field learning. However, wall-point velocities are constrained to zero by the CFD boundary conditions, providing no gradient information for the optimizer regardless of the loss weighting.

\subsubsection{Why Direct Prediction Succeeds}

Direct WSS prediction treats wall shear stress as an independent output variable, supervised directly by CFD ground truth at each wall point. This approach succeeds for several interconnected reasons.

First, each of the 145,000--388,000 wall points provides a meaningful WSS training signal, offering rich supervision rather than requiring the network to infer WSS from zero-velocity data. Second, the network can learn the complex spatial pattern of WSS independently from the velocity field through decoupled learning, allowing each output head to optimize for its specific task without interference. Third, the Navier-Stokes and continuity constraints provide physics regularization that ensures the velocity field remains physically consistent, even though WSS is not derived from it directly. The WSS physics constraint ($\mathcal{L}_{\text{wss-physics}}$) provides a soft coupling that encourages consistency between predicted WSS and velocity gradients without creating the gradient conflicts inherent in the derived approach.

\subsubsection{Implications for PINN Hemodynamics Research}

Our findings suggest that for PINN-based hemodynamic surrogate models trained on CFD data, direct prediction of surface quantities such as WSS and pressure with physics regularization outperforms approaches that attempt to derive these quantities from volume field predictions. This guidance is particularly relevant when training data enforces boundary conditions that create singularities, as with no-slip walls, when surface quantities are the primary clinical targets, and when high spatial resolution is required at boundaries. For inverse problems where the goal is to infer velocity fields from sparse measurements, the velocity-derived approach may still be appropriate. However, for surrogate modeling applications with abundant CFD training data, direct prediction is recommended.

\subsection{Physics-Informed Learning for Hemodynamics}

The physics-informed neural network approach demonstrated excellent performance in learning patient-specific WSS distributions while satisfying the governing Navier-Stokes equations. Several aspects contributed to this success.

The Fourier feature encoding proved essential for capturing high-frequency WSS patterns. Standard MLPs exhibited smoothing artifacts at bifurcations and stenoses where WSS gradients are steepest, but the random Fourier basis functions enable learning across a wider frequency spectrum. This capability is critical for hemodynamic applications where local flow features determine clinical outcomes.

Equally important was the mesh-based collocation sampling strategy. By sampling physics collocation points from the CFD mesh rather than the bounding box, we ensure that Navier-Stokes constraints are enforced only in physically relevant regions. This consideration is particularly important for coronary arteries, which are thin tubular structures occupying only a small fraction of their bounding volume.

The physics constraints also serve as powerful regularization during training. The Navier-Stokes and continuity losses enable accurate predictions while ensuring physically consistent velocity fields, distinguishing PINNs from purely data-driven approaches that may violate fundamental conservation laws.

\subsection{Clinical Applications}

The ability to predict WSS in real-time has significant clinical implications for CABG surgery across the entire perioperative timeline.

During pre-operative planning, surgeons can rapidly evaluate different graft configurations and anastomosis locations to optimize hemodynamic outcomes. The traditional approach requires multiple CFD simulations, each taking hours, which is impractical for clinical workflows where decisions must be made within constrained timeframes.

The real-time inference capability also opens possibilities for intra-operative guidance. WSS feedback could inform graft placement decisions during surgery itself, potentially improving long-term patency rates by avoiding configurations that create unfavorable hemodynamic conditions.

In the post-operative period, serial WSS assessments can track graft remodeling and identify early signs of failure without the burden of repeated CFD analyses. This capability enables proactive intervention before clinical symptoms manifest, potentially extending graft longevity and reducing the need for repeat revascularization procedures.

\subsection{Limitations and Future Work}

Several limitations of the current approach warrant discussion and suggest directions for future research.

The present models assume steady-state flow, neglecting the pulsatile nature of cardiac output. Extending the framework to time-dependent PINNs would capture important systolic and diastolic variations in WSS that may influence long-term vascular remodeling.

Our CFD training data and PINN physics losses employ the Newtonian fluid assumption. While appropriate for larger vessels where shear rates exceed 100 s$^{-1}$, blood exhibits non-Newtonian behavior at lower shear rates encountered in recirculation zones and near-wall regions. Incorporating the Carreau-Yasuda model into the PINN physics loss would improve accuracy in these clinically relevant low-flow regions.

The rigid wall assumption represents another simplification. Arterial compliance affects WSS through wall motion, and fluid-structure interaction extensions could address this limitation for more comprehensive hemodynamic characterization.

Finally, current models are trained separately for each patient, requiring 2--3 hours of computation per case. Developing transfer learning approaches that leverage pre-trained models could enable rapid adaptation to new patients with minimal fine-tuning, further accelerating clinical adoption.

% ============================================================================
% SECTION 5: CONCLUSION
% ============================================================================
\section{Conclusion}

This study presented a comprehensive computational investigation of blood flow in coronary arteries and saphenous vein grafts, combining traditional CFD analysis with physics-informed machine learning. CFD simulations using both Newtonian and Carreau-Yasuda blood models characterized WSS distributions across healthy, diseased, and post-CABG cases.

The FourierPINN surrogate models achieved excellent agreement with CFD simulations, with a mean NRMSE of 0.61\% and mean $R^2$ of 0.967 across 6 validated patient cases. The trained models predict WSS distributions at 145,000--388,000 points in under 1 second, representing a 3--4 orders of magnitude speedup over traditional CFD.

By embedding the incompressible Navier-Stokes equations directly into the learning process, the PINN ensures physically consistent predictions while maintaining computational efficiency for clinical translation. This approach bridges the gap between high-fidelity CFD simulations and clinical time constraints, enabling real-time hemodynamic assessment for CABG surgery planning and post-operative monitoring.

Future work will extend the framework to pulsatile flow, incorporate the Carreau-Yasuda non-Newtonian model in the PINN physics loss, and develop transfer learning approaches for rapid patient-specific adaptation.

% ============================================================================
% ACKNOWLEDGMENTS
% ============================================================================
\section*{Acknowledgments}

[Funding sources and acknowledgments]

% ============================================================================
% REFERENCES
% ============================================================================
\bibliographystyle{plain}
\begin{thebibliography}{99}

\bibitem{raissi2019physics}
Raissi, M., Perdikaris, P., \& Karniadakis, G. E. (2019). Physics-informed neural networks: A deep learning framework for solving forward and inverse problems involving nonlinear partial differential equations. \textit{Journal of Computational Physics}, 378, 686--707.

\bibitem{tancik2020fourier}
Tancik, M., Srinivasan, P., Mildenhall, B., et al. (2020). Fourier features let networks learn high frequency functions in low dimensional domains. \textit{Advances in Neural Information Processing Systems}, 33, 7537--7547.

\bibitem{spectral_bias}
Rahaman, N., Baratin, A., Arpit, D., et al. (2019). On the spectral bias of neural networks. \textit{International Conference on Machine Learning}, 5301--5310.

\bibitem{cabg_outcomes}
Gaudino, M., Benedetto, U., Bakaeen, F., et al. (2018). Off-pump versus on-pump coronary surgery and stroke. \textit{Journal of the American College of Cardiology}, 72(24), 3261--3270.

\bibitem{wss_graft_failure}
Owens, C. D. (2010). Adaptive changes in autogenous vein grafts for arterial reconstruction: clinical implications. \textit{Journal of Vascular Surgery}, 51(3), 736--746.

\bibitem{wss_thresholds}
Malek, A. M., Alper, S. L., \& Izumo, S. (1999). Hemodynamic shear stress and its role in atherosclerosis. \textit{JAMA}, 282(21), 2035--2042.

\bibitem{cfd_coronary}
Taylor, C. A., \& Figueroa, C. A. (2009). Patient-specific modeling of cardiovascular mechanics. \textit{Annual Review of Biomedical Engineering}, 11, 109--134.

\end{thebibliography}

\end{document}
